\subsection{模类集合的系数}

设 C 是 A-模的类的一个遗传集（VIII, 第 183 页，定义 1）。我们假设每个 C 型 A-模都在 K 上有限维，并用 \(\Theta_{\mathcal{C}}(\mathrm{A})\) 表示 C 型 A-模的系数所成的集合。由 VIII, 第 378 页的命题 2，集合 \(\Theta_{\mathcal{C}}(\mathrm{A})\) 也是诸 A-线性映射 \(u: E \to A^{*}\)（其中 E 取遍 C）的像的并；它也是子空间 \(\Theta_{\mathrm{E}}(\mathrm{A})\)（属于 \(A^{*}\)，其中 E 取遍 C）的并（同前引文）。(A, A)-子双模族 \((\Theta_{\mathrm{E}}(\mathrm{A}))_{\mathrm{E} \in \mathcal{C}}\)（\(A^{*}\) 的 (A, A)-子双模族）是有向的，故其并 \(\Theta_{\mathcal{C}}(\mathrm{A})\) 是 \(A^{*}\) 的一个 (A, A)-子双模。

命题 7. —— \(A^{*}\) 的一个 K 上有限维的左 A-子模包含于 \(\Theta_{\mathcal{C}}(\mathrm{A})\)，当且仅当它是 C 型的。

设 V 是 \(A^{*}\) 的一个 K 上有限维的左 A-子模。我们在 VIII, 第 379 页注 3 中已看到有 \(V \subset \Theta_{\mathrm{V}}(\mathrm{A})\)。若 V 是 C 型的，则 \(\Theta_{\mathrm{V}}(\mathrm{A}) \subset \Theta_{\mathcal{C}}(\mathrm{A})\)，故 V 包含于 \(\Theta_{\mathcal{C}}(\mathrm{A})\)。反之，设 V 包含于 \(\Theta_{\mathcal{C}}(\mathrm{A})\)。由于 \(\Theta_{\mathcal{C}}(\mathrm{A})\) 是有向族 \((\Theta_{\mathrm{E}}(\mathrm{A}))_{\mathrm{E}\in\mathcal{C}}\) 的并，且 V 在 K 上有限维，故存在一个 C 型模 E 使 V 包含于 \(\Theta_{\mathrm{E}}(\mathrm{A})\)。而存在自然数 n 使 \(\Theta_{\mathrm{E}}(\mathrm{A})\) 同构于 \(\mathrm{E}^{n}\) 的一个商模（VIII, 第 378 页，命题 2）。由于 C 是遗传的，V 是 C 型的。

推论. —— 设 E 是 K 上有限维的 A-模。那么 E 是 C 型的，当且仅当 \(\Theta_{\mathrm{E}}(\mathrm{A})\) 包含于 \(\Theta_{\mathcal{C}}(\mathrm{A})\)。

该条件显然是必要的。

反之，设 \(\Theta_{\mathcal{C}}(\mathrm{A})\) 包含 \(\Theta_{\mathrm{E}}(\mathrm{A})\)。A-模 \(\Theta_{\mathrm{E}}(\mathrm{A})\) 在 K 上有限维（VIII, 第 378 页，命题 2），故由命题 7 它是 C 型的。而存在整数 \(n \geqslant 0\) 使 E 同构于 \(\Theta_{\mathrm{E}}(\mathrm{A})^{n}\) 的一个 A-子模（VIII, 第 378 页，命题 2）。由于 C 是遗传的，E 是 C 型的。

